\documentclass{amsart}
\usepackage{amssymb}
\usepackage{enumitem}

\newtheorem{theorem}{Theorem}[section]
\newtheorem{lemma}[theorem]{Lemma}
\theoremstyle{definition}
\newtheorem{definition}[theorem]{Definition}
\theoremstyle{remark}
\newtheorem{remark}[theorem]{Remark}

\begin{document}

\title{A Note on Sums of Two Squares in Arithmetic Progressions}
\author{Helena Costa}
\address{Instituto de Matemática, Universidade Federal do Rio de Janeiro, Brazil}
\email{hcosta@example.br}
\author{Pieter de Vries}
\address{Mathematical Institute, Leiden University, The Netherlands}
\email{p.devries@example.nl}
\subjclass[2020]{Primary 11E25; Secondary 11N13}
\keywords{sums of two squares, arithmetic progressions}
\date{\today}

\begin{abstract}
We give an elementary proof that every arithmetic progression $a \bmod q$ with
$\gcd(a,q)=1$ and $a\equiv 1 \pmod 4$ contains infinitely many integers that are sums of
two squares, and we bound the least such integer.
\end{abstract}

\maketitle

\section{Introduction}
Let $\mathcal{S}$ denote the set of positive integers representable as $x^2+y^2$ with
$x,y\in\mathbb{Z}$. By Fermat's theorem, a prime $p$ lies in $\mathcal{S}$ if and only if
$p=2$ or $p\equiv1\pmod 4$.

\begin{theorem}\label{thm:main}
Let $q\ge1$ and $a$ be integers with $\gcd(a,q)=1$ and $a\equiv1\pmod 4$. Then the
progression $a\bmod q$ contains an element of $\mathcal{S}$ less than $q^{2}$.
\end{theorem}

\section{Preliminaries}
\begin{definition}
For $n\in\mathbb{N}$, let $r_2(n) = \#\{(x,y)\in\mathbb{Z}^2 : x^2+y^2=n\}$.
\end{definition}

\begin{lemma}\label{lem:mult}
The set $\mathcal{S}$ is closed under multiplication.
\end{lemma}

\begin{proof}
This follows from the Brahmagupta--Fibonacci identity
\[
(a^2+b^2)(c^2+d^2) = (ac-bd)^2 + (ad+bc)^2. \qedhere
\]
\end{proof}

\section{Proof of Theorem~\ref{thm:main}}
We proceed in three steps.
\begin{enumerate}[label=\textup{(\roman*)}]
  \item Reduce to the case $q$ odd using Lemma~\ref{lem:mult}.
  \item Choose $x$ with $x^2\equiv a-1\pmod q$ when possible.
  \item Otherwise use the Chinese remainder theorem on prime power factors of $q$.
\end{enumerate}

\begin{remark}
The exponent $2$ in Theorem~\ref{thm:main} is not optimal; Linnik-type arguments give
much smaller bounds.
\end{remark}

\bibliographystyle{amsplain}
\begin{thebibliography}{9}
\bibitem{HW} G.~H. Hardy and E.~M. Wright, \emph{An introduction to the theory of
numbers}, 6th ed., Oxford University Press, 2008.
\bibitem{IK} H.~Iwaniec and E.~Kowalski, \emph{Analytic number theory}, American
Mathematical Society, 2004.
\end{thebibliography}

\end{document}
